\chapter{Urine Crystals}

\section*{What Is This Test?}
A urine crystals test identifies and classifies crystalline structures in the urinary sediment. Crystals form when solutes (calcium, oxalate, uric acid, cystine, phosphate, struvite, drugs) exceed their solubility product in urine, a process influenced by urine pH, concentration, temperature, and the presence of inhibitors. Identifying crystals helps diagnose kidney stone risk, drug toxicity (methotrexate, sulfonamides, indinavir), and inherited metabolic disorders (cystinuria, primary hyperoxaluria). Crystals are detected by routine microscopic examination of the urine sediment, supplemented when needed by polarized light.

\section*{Alternative Names}
Urinary Crystals; Urine Microscopic Crystals; Crystalluria

\section*{Principle}
Crystals are visualized by light microscopy of fresh centrifuged urine at 400x. Some crystals (cholesterol, cystine) are best seen under polarized light. Crystal morphology is pH-dependent, so reading the dipstick pH first is essential. Some clinics use automated urine sediment analyzers (IQ200, iQ Cell, Sysmex UF-5000, Sedimax) using image recognition. Confirmation by chemical analysis or stone composition analysis may be necessary for atypical forms.

\section*{History}
Crystalluria was observed by early microscopists in the 17th--18th centuries. Clinical interest grew in the 20th century with the recognition of specific crystal-stone relationships (calcium oxalate, uric acid, struvite). Drug crystals (sulfonamides, antiretrovirals) gained prominence with modern pharmacotherapy.

\section*{Why This Test Is Ordered}
\begin{itemize}
  \item Evaluation of patients with hematuria, kidney stone passage, or flank pain
  \item Monitoring patients at high risk for urolithiasis (recurrent stone formers, inflammatory bowel disease, prior bariatric surgery, gout)
  \item Suspected drug-induced crystal nephropathy (methotrexate, sulfonamides, indinavir, acyclovir, triamterene)
  \item Evaluation of metabolic disorders (cystinuria, hyperoxaluria, gout)
  \item Routine urinalysis screening
\end{itemize}

\section*{Is This a Primary or Secondary Test}
Secondary test; part of routine urinalysis. Some crystals are pathologic and warrant further investigation.

\section*{How This Test Is Performed}
A fresh midstream urine sample (10--15 mL) is collected in a clean container. The specimen is centrifuged, supernatant decanted, resuspended in residual fluid, and examined microscopically at 400x. pH is recorded from dipstick. Crystals per high-power field are scored $0$--$4+$, described with morphology, and identified by pH and birefringence.

\section*{What Preparation Is Needed}
None for routine testing; for metabolic workup, a 24-hour urine collection may be requested. Avoid contrast media within 48 hours.

\section*{Contraindications and Precautions}
Crystals form as urine cools; refrigerated urine shows more crystals but is not clinically relevant if specimens are analyzed within 2 hours at room temperature. Alkalinization may dissolve uric acid crystals; acidification dissolves struvite. Some medication crystals may be confused with pathologic crystals; consult pharmacology references.

\section*{Risks}
None.

\section*{What Happens After the Test}
Report integrated with dipstick, specific gravity, and microscopy. Pathologic crystals warrant metabolic workup (24-hour urine, serum calcium, uric acid, cystine).

\section*{Understanding Results}

\subsection*{Below Normal}
Absence of crystals is normal and does not exclude stone risk; supersaturation dynamics are dynamic and intermittent.

\subsection*{Above Normal}
\begin{itemize}
  \item \textbf{Calcium oxalate:} Most common. Envelope or dumbbell shape; birefringent. Causes: hypercalciuria, hyperoxaluria, Crohn disease, high vitamin C, ethylene glycol poisoning
  \item \textbf{Calcium phosphate:} Alkaline urine; amorphous or stellate forms. Renal tubular acidosis, hyperparathyroidism
  \item \textbf{Uric acid:} Acidic urine (pH $<$ 5.5). Yellow-brown rhomboid shapes; birefringent. Gout, myeloproliferative disorders, chemotherapy tumor lysis
  \item \textbf{Struvite (magnesium ammonium phosphate):} Alkaline urine; coffin-lid shape. Urinary tract infection with urease-producing bacteria (Proteus, Klebsiella)
  \item \textbf{Cystine:} Hexagonal plates; pathognomonic for cystinuria (an autosomal recessive disorder)
  \item \textbf{Drug crystals:} Sulfonamide ``shocks of wheat''; methotrexate needles; indinavir starburst; acyclovir needle-shaped
\end{itemize}

\section*{Normal Results}
\begin{itemize}
  \item \textbf{No crystals observed or few crystals (0$-$1+):} Normal
  \item \textbf{Amorphous urates/phosphates:} Common; non-pathologic in small amounts
  \item \textbf{Crystals clinically significant:} Calcium oxalate in stone formers; uric acid in acidic urine with hyperuricosuria; cystine; struvite; drug crystals
  \item \textbf{Clinical context:} A few crystals in concentrated morning urine are often normal; their persistence in dilute urine is pathologic
\end{itemize}

\section*{Follow-up Tests}
\begin{itemize}
  \item \textbf{Stone analysis:} If patient has passed or surgery retrieved a stone, sent for infrared spectroscopy
  \item \textbf{24-hour urine metallolith study:} Calcium, oxalate, uric acid, cystine, citrate, sodium, volume
  \item \textbf{Serum calcium, PTH, uric acid:} Hyperparathyroidism, gout
  \item \textbf{Urine culture:} If struvite crystals to identify urease-producing bacteria
  \item \textbf{Renal ultrasound or noncontrast CT:} Confirm stones
  \item \textbf{Cystine (qualitative cyanide-nitroprusside test):} Detect cystinuria
\end{itemize}

\section*{References}
\begin{enumerate}
  \item Strasinger SK, Schumann GB. Urinalysis and Body Fluids. 7th ed. FA Davis; 2020.
  \item McPherson RA, Pincus MR, eds. Henry's Clinical Diagnosis and Management by Laboratory Methods. 24th ed. Elsevier; 2022.
  \item Pearle MS, Goldfarb DS, Assimos DG, et al. Medical management of kidney stones: AUA guideline. J Urol. 2014;192(2):316-324.
  \item Neumann I, et al. Drug-induced crystalluria. Am J Kidney Dis. 2020;76(1):134-142.
\end{enumerate}

\clearpage